\section{评估体系：为什么 ``evaluation is the key''}

训练信号、数据筛选和上线判断最终都依赖测量，因此评估不是后置考试，而是整个研发环路的控制器。本节在官方 slides 的 close/open-ended 区分之上，给出适合 Agent 的结果、过程、安全和成本四层口径，并说明自动判分与人类偏好各自承担什么角色。

\subsection{闭集评估与开集评估}
课程将评估分为 close-ended 与 open-ended。前者可以通过标准答案或结构化判定快速打分，后者通常没有唯一答案，需要多维标准与更复杂评审机制。

\begin{importantbox}{原话级观点}
讲者在约 01:23:20 强调：\textit{``evaluation is the key''}。这不是修辞，而是工程事实。没有高质量评估，训练改进无法可靠排序，资源投入会快速失效。
\end{importantbox}

\subsection{闭集评估的优缺点}
前一小节给出了两类评估的定义，这里先分析可自动验证的一侧。闭集任务适合高频回归和快速筛选，但固定答案空间也更容易被 prompt、解析器和数据泄漏影响；因此它提供的是低成本、低延迟的测量，不是对真实开放环境的完整替代。
闭集评估的优势是低成本、高自动化、易比较；缺点是可能被 prompt 和打分器偏差 ``刷分''。在 Agent 场景，闭集评估适合做快速回归，但不能单独作为上线依据。

\begin{knowledgebox}{闭集评估的推荐用途}
\begin{itemize}
    \item 每日训练回归与 smoke test；
    \item 新策略的快速筛选；
    \item 失败类型的早期预警。
\end{itemize}
\end{knowledgebox}

\subsection{开集评估的核心困难}
开集任务（如复杂客服、多轮协作、开放式写作）通常存在多解、长路径和偏好冲突。课程提醒，单一准确率指标不再有效，必须引入过程指标与结果指标共同约束。

\begin{warningbox}{开集评估中最常见的两个坑}
\begin{enumerate}
    \item 用 ``像人类'' 取代 ``完成目标''：观感好但任务失败。
    \item 用单次胜负取代稳定性：一次成功掩盖高方差问题。
\end{enumerate}
\end{warningbox}

\subsection{评估设计模板}
\begin{table}[H]
\centering
\begin{tabular}{p{3.0cm}p{4.0cm}p{4.8cm}}
\toprule
\textbf{层级} & \textbf{指标} & \textbf{示例} \\
\midrule
结果层 & 任务完成率、成功成本 & Ticket 关闭率、每单 token 花费 \\
过程层 & 错误恢复率、重计划次数 & API 失败后是否能自动纠错 \\
安全层 & 越权率、敏感操作拦截率 & 是否触发越权写入或危险调用 \\
体验层 & 延迟、可解释性评分 & 用户等待时间、行动理由可读性 \\
\bottomrule
\end{tabular}
\caption{课程观点落地为可执行评估模板}
\end{table}

\subsection{本章小结}
评估决定训练方向与资源分配，是 Agent 研发的 ``方向盘''。闭集评估负责速度，开集评估负责真实性，两者缺一不可。


